% !TEX root = ../algebra_02.tex

\section{Теорема Лагранжа и следствия из неё}

\begin{Theorem}{Теорема Лагранжа}{}
	Если $G$ --- конечная группа ($|G| < \infty$), а $H$ --- её подгруппа ($H < G$), то порядок группы равен произведению порядка подгруппы на её индекс:
	\[
	|G| = |H|\cdot (G:H)
	\].
\end{Theorem}

\begin{proof}
	Так как левые смежные классы образуют разбиение группы $G$ на непересекающиеся подмножества, порядок группы равен сумме мощностей всех смежных классов:
	\[
	|G| = \sum_{P \in G \slash H } |P|
	\].
	Рассмотрим произвольный смежный класс $P = gH$. Отображение $\eps \colon H \to gH$, заданное как $h \mapsto gh$, является биекцией (оно сюръективно по определению смежного класса и инъективно в силу закона сокращения в группе).
	Следовательно, мощность каждого смежного класса равна мощности подгруппы: $|gH| = |H|$.
	Отсюда:
	\[
	|G| = \sum_{P \in G \slash H} |H| = |H| \cdot |G \slash H| = |H| \cdot (G:H)
	\].
\end{proof}

\textbf{Следствия из теоремы Лагранжа:}
\begin{enumerate}
	\item Порядок любой подгруппы $H$ делит порядок конечной группы $G$: $|H| \mid |G|$.
	\item Порядок любого элемента $g \in G$ делит порядок группы: $\operatorname{ord}(g) \mid |G|$ (так как $\operatorname{ord}(g) = |\langle g \rangle|$).
	\item Любой элемент конечной группы, возведенный в степень, равную порядку группы, дает нейтральный элемент: $g^{|G|} = e$.
	\item Любая группа простого порядка $p$ является циклической. Действительно, если выбрать любой неединичный элемент $g \ne e$, его порядок $\operatorname{ord}(g)$ должен делить простое число $p$. Так как $\operatorname{ord}(g) \ne 1$, остается единственный вариант $\operatorname{ord}(g) = p = |G|$. Значит, группа порождается одним элементом.
\end{enumerate}
